%; whizzy paragraph -pdf xpdf -latex ./whizzypdfptex.sh
%; whizzy-paragraph "^\\\\begin{frame}\\|\\\\emtext"
% latex beamer presentation.
% platex, latex-beamer でコンパイルすることを想定。 

%     Tokyo Debian Meeting resources
%     Copyright (C) 2012 Junichi Uekawa

%     This program is free software; you can redistribute it and/or modify
%     it under the terms of the GNU General Public License as published by
%     the Free Software Foundation; either version 2 of the License, or
%     (at your option) any later version.

%     This program is distributed in the hope that it will be useful,
%     but WITHOUT ANY WARRANTY; without even the implied warreanty of
%     MERCHANTABILITY or FITNESS FOR A PARTICULAR PURPOSE.  See the
%     GNU General Public License for more details.

%     You should have received a copy of the GNU General Public License
%     along with this program; if not, write to the Free Software
%     Foundation, Inc., 51 Franklin St, Fifth Floor, Boston, MA  02110-1301 USA

\documentclass[cjk,dvipdfmx,12pt]{beamer}
\usetheme{Tokyo}
\usepackage{monthlypresentation}

%  preview (shell-command (concat "evince " (replace-regexp-in-string "tex$" "pdf"(buffer-file-name)) "&")) 
%  presentation (shell-command (concat "xpdf -fullscreen " (replace-regexp-in-string "tex$" "pdf"(buffer-file-name)) "&"))
%  presentation (shell-command (concat "evince " (replace-regexp-in-string "tex$" "pdf"(buffer-file-name)) "&"))

%http://www.naney.org/diki/dk/hyperref.html
%日本語EUC系環境の時
\AtBeginDvi{\special{pdf:tounicode EUC-UCS2}}
%シフトJIS系環境の時
%\AtBeginDvi{\special{pdf:tounicode 90ms-RKSJ-UCS2}}

\newenvironment{commandlinesmall}%
{\VerbatimEnvironment
  \begin{Sbox}\begin{minipage}{1.0\hsize}\begin{fontsize}{8}{8} \begin{BVerbatim}}%
{\end{BVerbatim}\end{fontsize}\end{minipage}\end{Sbox}
  \setlength{\fboxsep}{8pt}
% start on a new paragraph

\vspace{6pt}% skip before
\fcolorbox{dancerdarkblue}{dancerlightblue}{\TheSbox}

\vspace{6pt}% skip after
}
%end of commandlinesmall

\title{東京エリアDebian勉強会 with 第2回 Debianパッケージング道場}
\subtitle{第115回 2014年7月度}
\author{野島貴英}
\date{2014年7月19日}
\logo{\includegraphics[width=8cm]{image200607/openlogo-light.eps}}

\begin{document}

\begin{frame}
\titlepage{}
\end{frame}

\begin{frame}{設営準備にご協力ください。}
会場設営よろしくおねがいします。
\end{frame}

\begin{frame}{Agenda}
 \begin{minipage}[t]{0.45\hsize}
  \begin{itemize}
   \item 注意事項
	 \begin{itemize}
	  \item 写真はセミナールーム内のみ可です。
          \item 出入りは自由でないので、もし外出したい方は、まえださんまで一声くださいませ。
	 \end{itemize}
   \item 事前課題発表
   \item 第２回パッケージング道場
  \end{itemize}
 \end{minipage} 
 \begin{minipage}[t]{0.45\hsize}
  \begin{itemize}
   \item 最近あったDebian関連のイベント報告
	 \begin{itemize}
	  \item 第114回 東京エリアDebian勉強会
	 \end{itemize}
   \item Debian Trivia Quiz
   \item JenkinsでのDebian パッケージ 自動化
   \item 今後のイベント
   \item 今日の宴会場所
  \end{itemize}
 \end{minipage}
\end{frame}

\section{事前課題}
\emtext{事前課題}
{\footnotesize
 \begin{prework}{ koedoyoshida }
未定
\end{prework}

\begin{prework}{ dictoss(杉本　典充) }
パッケージ作成に関する作業は、\url{http://wnpp.debian.net/}に報告が出ているパッケージ化、またはパッチつくりをやろうと思います。具体的にはmunin、pygtkplot、pylzma、wx3.0-doc、hex2binが候補です。
\end{prework}

\begin{prework}{ 工藤寿晴 }
自分がVineで管理しているパッケージから、何かDebian用のパッケージを作ります。
（フォントかCUIのソフトウェアを検討中）
\end{prework}

\begin{prework}{ 吉野(yy\_y\_ja\_jp) }

manpages-ja
\end{prework}

\begin{prework}{ zinrai }

パッケージ作成

pulledpork\\
\url{https://code.google.com/p/pulledpork/}
\end{prework}

\begin{prework}{ wbcchsyn }

GitHub で公開中の自作プログラムをパッケージする予定
(先月の勉強会で発表した物です。)\\
\url{https://github.com/wbcchsyn/dVault}

時間があまったら、プログラムの機能実装を行います。
\end{prework}

\begin{prework}{ 鈴木崇文 }

Cocotのパッケージングをやります。
\end{prework}

\begin{prework}{ 野島　貴英 }

パッケージング道場参加します。
\begin{itemize}
\item xmrisのパッケージング
\item kanatestのバグつぶし
\item Softetherのパッケージ化ヒントが得られるまで 
\end{itemize}
となってます。
\end{prework}

\begin{prework}{ 岩松 信洋 }

パッケージアップロードします。
\end{prework}

\begin{prework}{ まえだこうへい }
\begin{itemize}
\item ToDo潰し\\
\url{http://d.palmtb.net/pages/about.html#todo-from-maintainer-dashboard}
\item pendingのままのSerf関連のパッケージ化
\end{itemize}
を行う予定。

\end{prework}

}

\section{第2回 Debianパッケージング道場}
\emtext{第2回 Debianパッケージング道場}

\begin{frame}[containsverbatim]{説明資料について(7/22野島追記)}
　今回のDebian パッケージング道場で、説明に利用されたプレゼン資料は、
「Debian パッケージングチュートリアル」となります。
\begin{commandlinesmall}
入手方法及び閲覧方法
$ sudo apt-get install packaging-tutorial
$ evince /usr/share/doc/packaging-tutorial/packaging-tutorial.ja.pdf
\end{commandlinesmall}
\end{frame}

\section{イベント報告}
\emtext{イベント報告}

\begin{frame}{第114回 東京エリアDebian勉強会}

 東京エリアDebian勉強会113回目は(株)スクウェア・エニックスさんで開催されました。
9名の参加者がありました。

 wbcchsynさんにて、GPGの秘密鍵を堅牢にかつセキュアに
管理する事を検討した成果についての発表がありました。今回は、その成果の１つとして、
堅牢化についてフォーカスし、RAIDの考え方を用いて複数の記憶メディアにデータを
分散配置して復元するソフトウェアをpythonで実装した事についての発表でした。
こちらについて様々に意見交換が参加者の間で行われました。

\end{frame}

\begin{frame}{第114回 東京エリアDebian勉強会}

　絶賛求む！！:発表者自身のアイデアに基づくDebianに関係した成果物の発表

\end{frame}


\section{Debian Trivia Quiz}
\emtext{Debian Trivia Quiz}
\begin{frame}{Debian Trivia Quiz}

  Debian の常識、もちろん知ってますよね?
知らないなんて恥ずかしくて、知らないとは言えないあんなことやこんなこと、
みんなで確認してみましょう。

今回の出題範囲は\url{debian-devel-announce@lists.debian.org},
\url{debian-devel@lists.debian.org} に投稿された
内容などからです。

\end{frame}

\subsection{問題}

%; whizzy-master ../debianmeetingresume201311.tex
% 以上の設定をしているため、このファイルで M-x whizzytex すると、whizzytexが利用できます。
%

\santaku
{libcが変更になりました。変更前と変更後の組み合わせで正しいのはどれ？}
{eglibcからglibcへ変更}
{glibcからeglibcへ変更}
{uClibcからglibcへ変更}
{A}
{eglibcが採択されてから5年が経過し、glibcに戻ってきた状況です。}

\santaku
{ネットワークはTorを使い、ファイル/メール/メッセンジャーソフトは暗号化対応という特徴を持つdebianベースのLive OSの名前は何でしょう？}
{Plan9}
{Tails}
{GNU/Hurd}
{B}
{TrailsはThe Amnesic Incognito Live Systemの略です。利用者の出来る限りの匿名性を保つ事を目的として作られたDebianベースのlive osとなります。}

\santaku
{Debian 6 LTS開発の為の寄付を受け付ける会社名は？}
{Canonical}
{SquareEnix}
{Freexian}
{C}
{Freexian社は、Debianの公式開発者のRapha\"el Hertzogにより設立された、
Debianの有償サポートをするための会社です。}

\santaku
{jessieのデフォルトのgccのバージョンは何になった？}
{4.9}
{4.6}
{2.6}
{A}
{各パッケージはgcc-4.9でもコンパイルが通るように修正が必要となります。}

\santaku
{jessieのfreezeは11月5日です（注：リリースじゃないよ。）ここで、スケジュールで新規のtransitionが出来なくなる日はいつでしょう？}
{10月5日}
{9月5日}
{8月5日}
{B}
{9月5日以降は新規のtransitionが出来なくなります。新しいパッケージをjessieに合わせてリリースしようと目論む方々は期日にご注意下さい。}






\section{Jenkins での Debian パッケージ自動化}
\emtext{Jenkins での Debian パッケージ自動化}

\section{今後のイベント}
\emtext{今後のイベント}
\begin{frame}{今後のイベント}
\begin{itemize}
 \item 関西エリアDebian勉強会
 \item 2014年8月 東京エリアDebian勉強会
\end{itemize}
\end{frame}

\section{今日の宴会場所}
\emtext{今日の宴会場所}
\begin{frame}{今日の宴会場所}
未定
\end{frame}

\end{document}

;;; Local Variables: ***
;;; outline-regexp: "\\([ 	]*\\\\\\(documentstyle\\|documentclass\\|emtext\\|section\\|begin{frame}\\)\\*?[ 	]*[[{]\\|[]+\\)" ***
;;; End: ***
